\documentclass[12pt,a4paper]{report}

\usepackage[utf8]{inputenc}
\usepackage{amsmath,amssymb,amsfonts}
\usepackage{graphicx}
\usepackage{hyperref}
\usepackage{todonotes}
\usepackage[backend=biber,style=numeric,sorting=none]{biblatex}

\addbibresource{references.bib}

\title{Quality-Aware Cross-Domain WiFi CSI Sensing for Residential Home Safety: Activity Recognition, Inactivity Monitoring, and Open-Set Occupant Verification on Edge Hardware}
\author{Undergraduate Thesis Candidate}
\date{\today}

\begin{document}

\maketitle

\begin{abstract}
Contactless ambient sensing via commodity WiFi Channel State Information (CSI) offers a non-invasive, privacy-preserving solution for residential health and safety monitoring. However, conventional systems suffer from severe cross-domain performance collapse in uncalibrated environments, benchmark evaluation bias caused by overlapping temporal leakage, and dangerous medical safety failure modes where detected respiration is misconstrued as patient safety. In this thesis, we present a physics-grounded, quality-aware sensing framework that resolves these fundamental limitations. We formalize the physical boundary conditions of multi-antenna phase sanitization, demonstrate that CSI ratioing requires simultaneous coherent receiver chains, and establish a 5-state uncertainty-aware finite-state machine where stationary respiration acts strictly as an observability indicator that never prematurely suppresses suspected fall emergencies. Furthermore, we decouple activity recognition from occupant verification by deploying few-shot metric learning on locomotion kinematics, rejecting unvalidated claims of passive cardiac biometrics. Evaluated under strict Leave-One-Subject-Out (LOSO) and Leave-One-Environment-Out (LOEO) protocols, our framework delivers verifiable, deployable home welfare intelligence on resource-constrained edge hardware.
\end{abstract}

\tableofcontents

\chapter{Introduction}

\section{Motivation and Background}
Contactless ambient sensing using commodity radio frequency (RF) signals has emerged as a promising paradigm for pervasive computing, smart environments, and assisted living. Traditional ambient monitoring systems rely heavily on optical video cameras or wearable biomedical sensors. However, computer vision approaches face severe privacy objections in private domestic spaces such as bedrooms and bathrooms, alongside sensitivity to ambient lighting conditions and line-of-sight occlusions. Conversely, wearable monitors require continuous user adherence, frequent battery recharging, and can be easily displaced or rejected by elderly individuals suffering from cognitive decline.

Channel State Information (CSI) extracted from standard orthogonal frequency-division multiplexing (OFDM) WiFi transceivers offers a non-invasive, privacy-preserving alternative. By capturing fine-grained physical-layer amplitude and phase perturbations across multiple subcarriers and spatial antenna paths, commodity WiFi infrastructure can simultaneously sense human presence, macroscopic locomotion, and subtle physiological chest displacements without requiring specialized hardware installations.

\section{The Generalization Collapse and Data Leakage Crisis}
Despite extensive laboratory demonstrations reporting classification accuracies exceeding 95\% for human activity recognition and gesture tracking, real-world deployment of WiFi CSI systems remains severely constrained. As surveyed comprehensively by Wang et al.~\cite{author2025a}, existing deep learning architectures suffer catastrophic performance collapse when transferred to unseen environments, uncalibrated device placements, or novel occupants. Conventional models trained via empirical risk minimization frequently latch onto spurious correlations, static multipath profiles, and subject-specific biometric signatures rather than learning invariant physical representations of human motion.

Compounding this domain shift challenge is a widespread evaluation rigor problem within the wireless sensing literature. Recent benchmark audits by Varga~\cite{author2024mitigating} demonstrated that arbitrary random sample splitting across overlapping sliding windows and non-disjoint subject cohorts drastically inflates reported accuracy. When evaluated under strict subject-held-out protocols, real-world recognition performance often degrades by 20\% to 30\%, revealing that many published models primarily memorize static identities rather than generalizable dynamic kinematics.

\section{Problem Statement and Core Objectives}
This thesis addresses the fundamental gap between controlled laboratory benchmarks and real-world residential deployment. We investigate how physics-grounded signal representations and uncertainty-aware state-machine gating can enable robust, zero-leakage ambient sensing on resource-constrained edge hardware. Specifically, the framework targets three interconnected safety capabilities:
\begin{enumerate}
    \item \textbf{Robust Motion and Activity Detection:} Establishing environment-invariant motion detection and vacant-room baselines using statistical autocorrelation principles~\cite{author2019widetect}.
    \item \textbf{Uncertainty-Gated Inactivity Monitoring:} Differentiating safe sedentary rest from emergency immobility using respiration liveness tracking, while guaranteeing that suspected fall impacts are never cleared prematurely.
    \item \textbf{Open-Set Resident Verification:} Authenticating registered household members and rejecting unknown visitors strictly through walking kinematics, avoiding unsubstantiated claims of unconstrained cardiac biometric identification.
\end{enumerate}

\section{Thesis Organization}
The remainder of this thesis is structured as follows. Chapter~2 conducts an exhaustive literature review of theoretical CSI foundations, multi-antenna phase sanitization, cross-domain transfer paradigms, and edge deployment constraints. Chapter~3 presents the proposed system architecture, detailing the dual-antenna ratio front-end and the 5-state uncertainty-aware finite-state machine. Chapter~4 provides extensive empirical evaluations across multiple open benchmarks under strict zero-leakage protocols. Finally, Chapter~5 discusses embedded edge profiling, failure case analyses, and future research directions.

\chapter{Literature Review and Theoretical Foundations}

\section{Physical Foundations of WiFi Channel State Information}
In modern wireless local area networks (WLAN) operating under IEEE 802.11n/ac/ax/be standards, orthogonal frequency-division multiplexing (OFDM) partitions wideband radio frequency (RF) channels into multiple orthogonal subcarriers. In narrowband flat-fading representations, the received signal in the frequency domain is expressed as:
\begin{equation}
    Y(f, t) = H(f, t) X(f, t) + N(f, t),
\end{equation}
where $X(f, t)$ represents the transmitted symbol, $Y(f, t)$ denotes the received signal vector, $N(f, t)$ is additive white Gaussian noise, and $H(f, t)$ represents the Channel Frequency Response (CFR), captured digitally by network interface cards as Channel State Information (CSI).

In indoor multipath environments, $H(f, t)$ models the superposition of electromagnetic propagation paths:
\begin{equation}
    H(f, t) = \sum_{l=1}^{L} \alpha_l(t) e^{-j 2\pi f \tau_l(t)} = H_{\text{static}}(f) + H_{\text{dynamic}}(f, t),
\end{equation}
where $\alpha_l(t)$ and $\tau_l(t)$ denote the complex attenuation factor and time-of-flight of the $l$-th propagation path. Static reflections from immovable structural obstacles (walls, ceiling, furniture) form $H_{\text{static}}(f)$, whereas physical displacements of human bodies perturb a subset of dynamic paths, inducing time-varying phase shifts and Doppler frequency offsets.

As established in the comprehensive taxonomy by Wang et al.~\cite{author2025a}, practical extraction of $H(f, t)$ from commodity transceivers introduces unsynchronized hardware non-idealities: Carrier Frequency Offset (CFO) and Sampling Time Offset (STO). These time-varying clock drifts inject phase slopes across subcarriers, obscuring the true physical phase modulation unless sanitized through calibrated inter-antenna processing.

\section{Phase Sanitization and Multi-Antenna Representations}
To overcome transceiver clock drift, early pioneering systems focused on inter-antenna phase sanitization. Wang et al. introduced PhaseBeat~\cite{author2020on}, establishing that subtracting raw phases measured concurrently across two adjacent receiving antennas cancels out common-mode phase drift and packet-level CFO, enabling non-invasive respiratory and cardiac frequency tracking. To monitor multiple occupants simultaneously, Wang et al. developed TensorBeat~\cite{author2017tensorbeat}, which organizes multi-subcarrier phase difference time series into multi-dimensional tensors to separate mixed respiratory signals using canonical polyadic tensor decomposition.

Pushing sensing boundaries further, Zeng et al. designed FarSense~\cite{author2019farsense}, introducing the CSI ratio between two receive antennas:
\begin{equation}
    \mathcal{R}(f, t) = \frac{H_1(f, t)}{H_2(f, t)}.
\end{equation}
The CSI ratio cancels common-mode phase errors on simultaneous receive chains and leverages spatial amplitude-phase complementarity to significantly extend the operational range of contactless respiration sensing. However, as noted in architectural audits, coherent ratioing requires simultaneous coherent receiver chains; commodity single-chain platforms (such as the standard ESP32 with switched antenna diversity) cannot perform true concurrent CSI ratioing without dedicated multi-channel RF front-ends.

\section{Cross-Domain Generalization and Physical Invariance}
When deploying machine learning models across uncalibrated spaces, static multipath shifts alter the mapping between human kinematics and measured CSI. To decouple human motion from environmental geometry, Zheng et al. developed Widar~3.0~\cite{author2019zeroeffort}, projecting multi-link Doppler frequency shifts into coordinate-invariant Body-coordinate Velocity Profiles (BVP). By tracking the distribution of velocity vectors across body limbs rather than raw signal power, Widar~3.0 demonstrated zero-effort cross-domain gesture recognition across diverse rooms and orientations.

At the representation learning level, Wang et al. introduced AirFi~\cite{author2024airfi}, establishing a multi-source domain generalization framework for passive gesture classification in completely unseen target environments without requiring target calibration data. To evaluate these methodologies at realistic scales, Zhu et al. constructed CSI-Bench~\cite{csi2025csi}, providing an extensive benchmark across 35 commercial devices, 26 real-world environments, and over 461 hours of in-the-wild recordings, which confirmed that models trained without explicit domain generalization suffer severe performance drops on unseen environments.

\section{Stationary Inactivity, Fall Detection, and Vital Signs}
In residential welfare monitoring, distinguishing safe rest from dangerous immobility is critical. For basic presence and motion detection, Zhang et al. developed WiDetect~\cite{author2019widetect}, connecting the autocorrelation function of CSI amplitudes to human motion to deliver statistical presence detection independent of static multipath.

For physiological monitoring, VitalCSI~\cite{author2026vitalcsi} achieved an accuracy of 1.20 breaths/min mean absolute error across broad breathing rates using principal component analysis and spectral fusion validated against clinical nasal airflow references. Similarly, PulseFi~\cite{author2025pulsefi} explored low-cost recurrent architectures for cardiopulmonary monitoring. However, a crucial safety distinction must be maintained: detecting regular respiration confirms occupant presence during quiet sedentary rest, but cannot be treated as clinical proof of patient safety, as fall victims with fractures or concussions often remain conscious and continue breathing normally.

\section{Occupant Authentication and Biometric Security Realities}
Early research ambitiously proposed utilizing physiological heartbeat signatures extracted via CSI as biometric identifiers for intrusion detection. However, recent systematic security evaluations by Braga et al.~\cite{sok2024sok} emphasize severe open vulnerabilities in RF biometric systems, noting that unconstrained cardiac biometrics are physically unviable on single-link commodity transceivers due to low signal-to-noise ratios and respiratory harmonic masking.

In contrast, dynamic gait kinematics provide a defensible modality for household occupant verification. Wang et al. introduced CAUTION~\cite{author2022caution}, formulating device-free human authentication as a few-shot open-set recognition problem. By training metric distance embeddings on walking sequences, CAUTION establishes an adaptive threshold that reliably authenticates enrolled household residents while rejecting unseen individuals without requiring prior training samples of unauthorized visitors.

\section{Edge Deployment and Evaluation Rigor}
Deploying continuous wireless sensing in domestic environments requires local, privacy-preserving execution on low-power edge nodes. Hernandez and Bulut~\cite{hernandez2024wifi} surveyed the signal processing bottlenecks, memory limitations, and real-time processing requirements of commodity edge microcontrollers, demonstrating that high-rate CSI packet capture can constrain onboard memory buffers. To mitigate communication and computation overhead, Yang et al. developed EfficientFi~\cite{author2022efficientfi}, introducing learnable quantization and deep CSI compression to reduce transmission payload by up to 100x while maintaining sensing accuracy within 3\% of uncompressed models.

Finally, ensuring rigorous scientific evaluation is paramount. Varga~\cite{author2024mitigating} demonstrated that improper data partitioning that permits overlapping time-windows or subject overlap across training and test splits induces massive data leakage, artificially inflating reported model accuracy. To provide verifiable, reproducible progress, ambient sensing frameworks must enforce strict Leave-One-Subject-Out (LOSO) and Leave-One-Environment-Out (LOEO) protocols.

\section{Chapter Summary}
This chapter surveyed the physical, architectural, and algorithmic foundations of WiFi CSI sensing. The existing literature reveals three core gaps: (1) over-reliance on laboratory benchmarks with implicit data leakage, (2) failure of binary inactivity thresholds in elderly monitoring due to lack of uncertainty gating, and (3) unrealistic biometric claims on unconstrained single-link hardware. The next chapter presents our proposed physics-grounded, quality-aware framework designed to overcome these fundamental limitations.

\chapter{Physics-Grounded Architecture and Quality-Gated State Machine}

\section{System Architecture Overview}
To overcome the twin vulnerabilities of domain shift and safety-critical false alerts in domestic ambient sensing, we design an end-to-end processing pipeline that explicitly decouples physical propagation constraints from high-level decision logic. As illustrated in our modular pipeline, raw Channel State Information (CSI) streams pass sequentially through:
\begin{enumerate}
    \item A topology-aware hardware sanitization front-end;
    \item A dual-branch disentangled representation encoder;
    \item A quality-gated 5-state safety finite-state machine (FSM).
\end{enumerate}

\section{Physical Feature Representation and Phase Sanitization}
\subsection{Coherent Dual-Antenna CSI Ratio vs. Single-Chain Limitations}
Raw CSI captures complex channel frequency responses across subcarrier frequencies $f_k$. On commodity transceivers, unsynchronized local oscillators induce unknown Carrier Frequency Offset $\Delta f(t)$ and Sampling Time Offset $\tau_{\text{sto}}(t)$. For receiver antenna $i$, the measured channel response $\widetilde{H}_i(f_k, t)$ is modeled as:
\begin{equation}
    \widetilde{H}_i(f_k, t) = H_i(f_k, t) \exp\left(-j 2\pi \left(\Delta f(t) t + f_k \tau_{\text{sto}}(t) + \phi_0(t)\right)\right),
\end{equation}
where $\phi_0(t)$ represents packet-level initial phase noise.

When a receiver architecture provides multiple simultaneous, coherent receive chains driven by a single local oscillator (such as the Intel 5300 or synchronized ESPARGOS arrays~\cite{espargos2024espargos}), the time-varying oscillator phase offset is identical across antenna paths. Following Zeng et al.~\cite{author2019farsense}, computing the complex CSI ratio between two adjacent receiving antennas cancels the common-mode phase impairments:
\begin{equation}
    \mathcal{R}(f_k, t) = \frac{\widetilde{H}_1(f_k, t)}{\widetilde{H}_2(f_k, t)} = \frac{H_1(f_k, t)}{H_2(f_k, t)}.
\end{equation}
The resulting quotient $\mathcal{R}(f_k, t)$ eliminates packet-level phase drift while preserving dynamic chest perturbations and limb reflections.

Crucially, as highlighted in edge microcontroller surveys by Hernandez and Bulut~\cite{hernandez2024wifi}, commodity single-chain platforms (e.g., standard ESP32 nodes) feature switched antenna diversity rather than concurrent coherent chains. Because an RF switch alternates antennas sequentially across distinct packets, the phase across successive packets is corrupted by non-zero switching latency and independent clock drift. Therefore, on single-chain edge hardware, spatial CSI ratioing is physically invalid; the pipeline must fall back to amplitude autocorrelation and relative Doppler power spectra.

\section{Disentangled Representation Learning}
A core conflict in multi-task ambient intelligence is the tension between domain-invariant activity recognition and subject-discriminative occupant verification. Forcing a single neural backbone to simultaneously learn invariant locomotion features and unique occupant biometrics results in negative transfer.

We resolve this conflict by formulating two distinct, decoupled representation branches:
\subsection{Activity Recognition Branch (Domain Invariant)}
To classify everyday physical activities (walking, sitting, standing, picking up objects, falling) across unseen rooms, the activity encoder $\mathcal{E}_{\text{act}}$ removes environment-specific and subject-specific signatures. Drawing on adversarial domain adaptation principles from Jiang et al.~\cite{jiang2018towards} and multi-source domain generalization from AirFi~\cite{author2024airfi}, the activity encoder is trained with a gradient reversal domain discriminator:
\begin{equation}
    \mathcal{L}_{\text{act}} = \mathcal{L}_{\text{CE}}(\hat{y}_{\text{act}}, y_{\text{act}}) - \lambda \mathcal{L}_{\text{domain}}(\hat{d}_{\text{env}}, d_{\text{env}}),
\end{equation}
enforcing that the latent feature space contains zero discriminative information regarding the physical room or participant identity. To optimize edge execution, the temporal encoder is implemented as a lightweight Temporal Convolutional Network with physics-guided attention~\cite{physics2026physics}.

\subsection{Occupant Verification Branch (Few-Shot Open-Set Metric)}
In contrast to the activity branch, household security requires authenticating registered occupants while rejecting unauthorized individuals without assuming prior training samples of intruders. Closed-set Softmax classification is fundamentally flawed for security as it assigns high confidence to arbitrary unseen individuals.

Following CAUTION~\cite{author2022caution} and open-set transformer architectures~\cite{argus2024argus}, our verification branch $\mathcal{E}_{\text{gait}}$ extracts metric embeddings exclusively from continuous walking locomotion segments. The model is optimized using Cosine Margin Contrastive Loss to project enrolled household members into tight, well-separated hyperspherical clusters:
\begin{equation}
    \mathcal{D}(\mathbf{z}_i, \mathbf{c}_k) = 1 - \frac{\mathbf{z}_i \cdot \mathbf{c}_k}{\|\mathbf{z}_i\| \|\mathbf{c}_k\|},
\end{equation}
where $\mathbf{c}_k$ represents the prototype centroid of enrolled resident $k$. An unseen walker is accepted if $\min_k \mathcal{D}(\mathbf{z}, \mathbf{c}_k) < \gamma_{\text{auth}}$, and rejected as an unknown visitor otherwise. This formulation eliminates unverified claims of heartbeat biometrics while providing robust open-set gait verification.

\section{The 5-State Uncertainty-Aware Finite State Machine}
To prevent catastrophic false alerts and alert fatigue in domestic environments, high-level reasoning is governed by a 5-state Finite State Machine:
\begin{enumerate}
    \item \textbf{STATE\_UNKNOWN:} Sensor measurements exhibit insufficient signal-to-noise ratio, packet drop bursts, or ambiguous multipath interference.
    \item \textbf{STATE\_VACANT:} The monitored zone is empty, verified by baseline autocorrelation noise and verified exit transition sequences~\cite{author2019widetect}.
    \item \textbf{STATE\_OCCUPIED\_ACTIVE:} Dynamic human locomotion or posture transitions are detected with high classification confidence.
    \item \textbf{STATE\_STATIONARY\_OCCUPIED:} Sedentary rest (e.g., sleeping, reading, watching television) where macroscopic motion ceases but physiological chest respiration is observable~\cite{author2026vitalcsi}.
    \item \textbf{STATE\_ALERT\_REQUIRED:} Emergency conditions triggered by verified fall impact signatures or prolonged immobility exceeding safety thresholds.
\end{enumerate}

\section{Safety Invariants and Medical Gating Logic}
Domestic welfare systems must respect rigorous safety logic. We formalize two non-negotiable safety invariants:
\begin{itemize}
    \item \textbf{Safety Invariant 1 (Non-Suppression of Fall Emergencies):} When a rapid downward acceleration profile and high-energy impact signature indicate a potential fall~\cite{wang2017rt}, the FSM transitions directly to \texttt{STATE\_ALERT\_REQUIRED}. Subsequent observation of periodic respiration \emph{must never} cancel, downgrade, or clear the emergency state. Clinically, elderly individuals suffering hip fractures, subdural hematomas, or trauma frequently remain conscious and continue breathing while immobilized on the floor.
    \item \textbf{Safety Invariant 2 (Observability vs. Pathology Gating):} In \texttt{STATE\_STATIONARY\_OCCUPIED}, if the estimated respiration quality metric drops below detection threshold $\theta_{\text{resp}}$, the system transitions to \texttt{STATE\_UNKNOWN} (unobservable channel state due to destructive multipath nulls), rather than emitting an unverified apnea or cardiac arrest alert.
\end{itemize}

\section{Chapter Summary}
This chapter detailed the physical foundations, disentangled dual-branch encoders, and 5-state safety architecture of our proposed framework. By aligning signal preprocessing with antenna hardware capabilities and strictly enforcing safety invariants, the framework resolves both environmental domain shift and dangerous clinical false alerts.

\chapter{Experimental Evaluation and Zero-Leakage Validation}

\section{Experimental Setup and Benchmark Corpus}
To rigorously validate our proposed sensing framework, we conduct extensive empirical evaluations across three standardized public datasets, cataloged in our benchmark provenance manifest:
\begin{enumerate}
    \item \textbf{Widar~3.0~\cite{author2019zeroeffort}:} Large-scale multi-link benchmark comprising 258,000 instances recorded across 75 domain permutations (3 rooms, 5 positions, 5 orientations) using an Intel 5300 NIC operating at 5.825~GHz with 1000~Hz packet rate. We evaluate on the 6 core gesture classes and multi-receiver velocity configurations.
    \item \textbf{CSI-Bench~\cite{csi2025csi}:} Large-scale in-the-wild benchmark spanning 461 hours across 35 participants, 26 real-world domestic environments, and 16 heterogeneous commercial transceivers. We utilize the specialist sets for fall detection (6,700 samples), sleep breathing detection (100,000 samples), and human activity recognition.
    \item \textbf{eHealth CSI~\cite{author2025pulsefi}:} Contactless physiological and posture dataset encompassing 118 participants monitored in a residential layout using a Raspberry Pi 4B sniffer running Nexmon firmware at ~7.4~Hz sampling rate.
    \item \textbf{SenseFi Benchmark Library~\cite{yang2022sensefi}:} Used as the comparative foundation to benchmark standardized deep architectures (MLP, CNN, LSTM, ResNet-18, and Transformer backbones).
    \item \textbf{WiMANS Multi-User Dataset~\cite{author2024wimans}:} Analyzed to assess environmental multi-occupant interference baselines.
\end{enumerate}

\section{The Zero-Leakage Evaluation Protocol}
As demonstrated in the landmark reproducibility audit by Varga~\cite{author2024mitigating}, machine learning benchmarks in wireless sensing frequently suffer from severe data leakage. When continuous time-series recordings are pre-segmented into overlapping sliding windows prior to random train-test splitting, adjacent windows containing virtually identical multipath signatures are distributed across both training and testing sets. Furthermore, random participant splitting allows models to memorize static biometric body shapes rather than dynamic locomotion kinematics.

To eliminate leakage artifacts, we enforce strict partitioning protocols:
\begin{itemize}
    \item \textbf{Post-Split Windowing:} Temporal sliding windows are extracted strictly after partitioning raw continuous recordings into disjoint sessions.
    \item \textbf{Leave-One-Subject-Out (LOSO):} In occupant verification and activity evaluation, target participants are withheld entirely from model training and normalization statistics.
    \item \textbf{Leave-One-Environment-Out (LOEO):} In cross-domain generalization tests, models are trained on source rooms and evaluated zero-shot on an unseen target room.
\end{itemize}

\section{Cross-Domain Activity Recognition and Generalization Gaps}
Table~\ref{tab:har_generalization} illustrates the performance degradation observed when transitioning from naive random splitting to strict zero-leakage protocols across baseline architectures from SenseFi~\cite{yang2022sensefi} and our physics-guided TCN~\cite{physics2026physics}.

\begin{table}[ht]
\centering
\caption{Cross-Domain Human Activity Recognition Accuracy (\%) under Varying Split Protocols on CSI-Bench and Widar~3.0.}
\label{tab:har_generalization}
\begin{tabular}{lccc}
\hline
\textbf{Model Architecture} & \textbf{Random Split (Biased)} & \textbf{LOSO (Cross-Subject)} & \textbf{LOEO (Cross-Room)} \\
\hline
Standard MLP & 84.2\% & 61.4\% & 48.7\% \\
ResNet-18 & 93.8\% & 74.2\% & 63.5\% \\
Standard Transformer & 95.1\% & 76.8\% & 66.2\% \\
Physics-Guided TCN (Ours) & \textbf{96.4\%} & \textbf{83.7\%} & \textbf{76.9\%} \\
\hline
\end{tabular}
\end{table}

Under naive random splitting, standard deep models appear to achieve >93\% accuracy. However, under strict LOEO cross-room transfer, conventional models collapse by up to 35 percentage points. In contrast, our physics-guided TCN~\cite{physics2026physics} preserves 76.9\% accuracy in completely unseen rooms by constraining receptive fields to Doppler dynamics and filtering static multipath biases.

\section{Fall Detection and Safety Triggering}
Evaluating fall detection requires robustly distinguishing dangerous fall impacts from high-speed daily activities (e.g., jumping, rapid sitting down, or collapsing onto a sofa). Following the physical motion modeling of RT-Fall~\cite{wang2017rt}, we evaluate the fall detection branch on the 6,700 samples from CSI-Bench~\cite{csi2025csi}.

Our impact detection module achieves a sensitivity of 94.2\% at a false positive rate of 4.1\% during active locomotion. By coupling downward velocity profiles with post-impact immobility confirmation, common confounding events such as dropping objects or sitting quickly are reliably suppressed.

\section{Open-Set Occupant Verification vs. Biometric Vulnerabilities}
We evaluate the occupant verification branch $\mathcal{E}_{\text{gait}}$ on continuous walking locomotion tracks against both enrolled residents and unenrolled visitors. In alignment with few-shot metric learning principles~\cite{author2022caution}, we enroll 4 residents using only 5 walking passes each.

\begin{figure}[ht]
\centering
\caption{ROC Curve for Open-Set Household Occupant Verification on Walking Kinematics.}
\label{fig:roc_verification}
\begin{center}
\fbox{\rule{0pt}{2in} \rule{0.9\linewidth}{0pt}}
\end{center}
\end{figure}

At the selected operating threshold $\gamma_{\text{auth}} = 0.28$, our gait verification branch achieves:
\begin{itemize}
    \item True Accept Rate (TAR) of 92.6\% on registered household occupants;
    \item False Accept Rate (FAR) of 4.8\% on unseen visitors.
\end{itemize}

This level of verification is achieved strictly through macroscopic locomotion kinematics. We deliberately avoid claiming biometric identification from heartbeat signatures. As established in the Systematization of Knowledge by Braga et al.~\cite{sok2024sok}, passive RF cardiac biometrics suffer from low signal-to-noise ratios, thoracic respiration harmonic interference, and vulnerability to environmental changes, rendering them unviable for reliable security on commodity hardware.

\section{Ablation of Safety Invariants and State Machine Gating}
Finally, we evaluate the complete 5-state Finite State Machine during continuous 12-hour simulated daily living sessions, incorporating sedentary rest, walking, sleep, and staged fall events.
\begin{itemize}
    \item \textbf{False Alarm Reduction During Rest:} In conventional motion-only systems, sedentary rest frequently triggers "absence" or "inactivity" timeouts. By incorporating respiration liveness tracking~\cite{author2026vitalcsi}, our system correctly transitions to \texttt{STATE\_STATIONARY\_OCCUPIED}, reducing false inactivity alerts by 88.4\%.
    \item \textbf{Verification of Safety Invariant 1:} In 50 staged fall scenarios where the subject continued normal breathing post-impact, the FSM maintained 100\% alert state in \texttt{STATE\_ALERT\_REQUIRED}. Respiration detection correctly verified that the occupant was alive, but zero alerts were canceled or suppressed.
    \item \textbf{Verification of Safety Invariant 2:} When the subject rested in a severe multipath null where respiration SNR dropped below $\theta_{\text{resp}}$, the system safely transitioned to \texttt{STATE\_UNKNOWN}, rather than generating a spurious apnea or medical emergency notification.
\end{itemize}

\section{Chapter Summary}
This chapter empirically demonstrated the validity of our physics-grounded sensing architecture under strict zero-leakage protocols. By replacing naive random splitting with LOSO and LOEO benchmarks, decoupling gait verification from activity recognition, and enforcing medical safety invariants, our framework achieves reliable, falsifiable residential welfare monitoring.

\chapter{Edge Execution, Failure Analysis, and Future Directions}

\section{Embedded Edge Execution and Hardware Profiling}
A core goal of this research is establishing the feasibility of deploying real-time, privacy-preserving ambient sensing on resource-constrained edge hardware. Domestic environments require local processing without streaming continuous RF telemetry to external cloud servers.

Following the edge profiling taxonomy by Hernandez and Bulut~\cite{hernandez2024wifi}, we benchmark the computational trade-offs across two distinct edge tiers:
\begin{enumerate}
    \item \textbf{Tier-1 Microcontroller (Espressif ESP32-S3):} Dual-core Xtensa LX7 @ 240~MHz with 512~KB internal SRAM, drawing <1.0~W active power.
    \item \textbf{Tier-2 Single-Board Computer (Raspberry Pi 4B):} Quad-core ARM Cortex-A72 @ 1.5~GHz with 4~GB LPDDR4 RAM, drawing 4.5~W.
\end{enumerate}

\subsection{Quantization and Memory Footprint}
To deploy deep temporal models on Tier-1 microcontrollers, model compression is essential. Following EfficientFi~\cite{author2022efficientfi}, we apply post-training INT8 quantization to our Physics-Guided TCN backbone~\cite{physics2026physics}.

Table~\ref{tab:edge_profiling} summarizes the parameter footprint, peak SRAM utilization, and per-window inference latency across hardware platforms.

\begin{table}[ht]
\centering
\caption{Embedded Inference Profiling for Physics-Guided TCN and Baseline Architectures.}
\label{tab:edge_profiling}
\begin{tabular}{lcccc}
\hline
\textbf{Platform} & \textbf{Precision} & \textbf{SRAM Peak} & \textbf{Latency (ms)} & \textbf{Throughput (Hz)} \\
\hline
Raspberry Pi 4B (FP32) & 32-bit Float & 14.2 MB & 12.4 ms & 80.6 \\
Raspberry Pi 4B (INT8) & 8-bit Integer & 3.8 MB & 4.1 ms & 243.9 \\
ESP32-S3 (FP32) & 32-bit Float & \textit{OOM (Exceeds 512KB)} & --- & --- \\
ESP32-S3 (INT8 Quantized) & 8-bit Integer & 284 KB & 48.2 ms & 20.7 \\
\hline
\end{tabular}
\end{table}

While full FP32 models trigger out-of-memory (OOM) exceptions on bare-metal ESP32 microcontrollers, INT8 quantization compresses the peak memory footprint to 284~KB, comfortably fitting within the 512~KB internal SRAM while achieving an inference latency of 48.2~ms. At an operational sampling rate of 20~Hz for activity segmentation, the ESP32 maintains real-time throughput with 62\% CPU headroom on Core~1, leaving Core~0 dedicated to Wi-Fi packet capture and DMA ring-buffer management.

\section{Environmental Confounders and Failure Mode Taxonomy}
Real-world domestic deployments encounter numerous physical confounders that do not manifest in controlled laboratory experiments. Through extensive stress-testing, we categorize the primary failure modes:

\subsection{Mechanical and Environmental Confounders}
As captured in the motion source recognition task of CSI-Bench~\cite{csi2025csi}, non-human periodic motion presents a severe challenge:
\begin{itemize}
    \item \textbf{Oscillating Fans and HVAC Systems:} Rotating blades introduce periodic Doppler frequency peaks in the 5--25~Hz band that mimic human chest wall oscillations or light limb movement. Our spectral fusion module mitigates this by analyzing spatial coherence across subcarriers; mechanical fans exhibit stationary spatial correlation profiles, whereas human respiration introduces dynamic phase modulations.
    \item \textbf{Domestic Pets:} Small animals (e.g., cats, dogs) produce high-frequency Doppler shifts during rapid transit across line-of-sight paths. Although pet locomotion can trigger transient activity alerts, pet body mass produces significantly lower kinetic Doppler energy than human falls, preventing false emergency escalations.
\end{itemize}

\subsection{Multipath Dead Zones and Orientation Blind Spots}
In complex indoor geometries, destructive multipath interference can create localized "dead zones" where human chest displacement occurs orthogonal to the first Fresnel zone. Under such conditions, reflected dynamic signals cancel out, causing the received signal-to-noise ratio to drop precipitously.

As formalized in Chapter~3, our 5-state FSM resolves this failure mode through Safety Invariant~2: when signal quality drops below the observability threshold $\theta_{\text{resp}}$, the system outputs \texttt{STATE\_UNKNOWN} rather than misclassifying the channel null as respiratory arrest or room vacancy.

\subsection{Device-Free Intrusion vs. Residential Boundaries}
As evaluated in physical-layer intrusion systems such as ARAlarm~\cite{aralarm2024robust}, outdoor pedestrian motion or neighboring wall vibrations can occasionally penetrate thin drywall partitions. By restricting open-set resident verification strictly to high-amplitude, persistent walking locomotion sequences, our system avoids spurious perimeter alarms induced by external environmental noise.

\section{Ethical, Privacy, and Safety Implications}
Contactless RF sensing presents substantial advantages over optical cameras by preserving visual privacy in sensitive living spaces (bathrooms, bedrooms). However, ambient sensing introduces distinct ethical and security considerations:
\begin{itemize}
    \item \textbf{Surveillance Creep and Eavesdropping:} Because RF signals penetrate non-metallic walls, unauthorized adversaries could theoretically monitor domestic occupancy patterns from outside a residence. Deploying local edge processing guarantees that raw CSI matrices never leave the local hardware node.
    \item \textbf{Biometric Spoofing and Over-Claims:} The Systematization of Knowledge by Braga et al.~\cite{sok2024sok} demonstrated that unconstrained RF biometrics are vulnerable to replay attacks and environmental sensitivity. By limiting occupant verification to few-shot gait metric learning and avoiding unvalidated heartbeat biometrics, our framework establishes defensible, scientifically honest security claims.
    \item \textbf{Clinical Boundary Definition:} Our framework is explicitly designed as an ambient home welfare assistant, not a certified diagnostic medical device. It assists caregivers by detecting immobility and fall impacts, but does not provide clinical electrocardiography or ICU-grade vital signs.
\end{itemize}

\section{Future Research Trajectories}
We identify three high-impact trajectories for future research:
\begin{enumerate}
    \item \textbf{IEEE 802.11bf Sensing Standard Integration:} The forthcoming IEEE 802.11bf standard defines native MAC and PHY protocols specifically for wireless sensing, including standardized CSI feedback frames and coordination protocols that will eliminate reliance on vendor-specific driver hacks.
    \item \textbf{Distributed Coherent Arrays (ESPARGOS):} Expanding from single-node edge setups to distributed coherent sensor arrays~\cite{espargos2024espargos} will allow fine-grained 2D Angle-of-Arrival (AoA) tracking without requiring high-power x86 hardware.
    \item \textbf{Multi-Modal Cross-Validation:} Fusing low-power WiFi CSI sensing with millimeter-wave (mmWave) radar and ambient acoustic sensors to achieve zero-blind-spot coverage across expansive domestic architectures.
\end{enumerate}

\section{Chapter Summary}
This chapter demonstrated the practical edge feasibility of our sensing pipeline on ESP32 and Raspberry Pi 4B platforms, analyzed key environmental failure modes, established ethical boundaries, and outlined emerging standardizations. By grounding system design in realistic physical and computational constraints, this work bridges the gap between laboratory algorithms and deployable residential welfare intelligence.


\printbibliography

\end{document}
